\subsection{A joint residual trace criterion from the source model}
\label{sec:joint-trace-source}

The conclusions in this section retain the stipulated arithmetic source
inputs of Lemma~\ref{lem:step17-unit}. They are deductions from the
complete intermediate proofs at the pinned commit
\cite{K17PinnedLeanMirror}; no new Lean compilation is asserted. Put
\[
 \mathcal A=\mathrm{Adm3U\_s}(\iota,e,N)
       \cup\mathrm{Adm2A}(\iota,e,N|d_K|),\qquad
 t(s)=\Psi(s)+1+\Psi(s)^{-1},\quad \Psi=\chi_\ell^e\psi,
\]
where the characters are restricted to $G_K$. Write
$\mathfrak m_{\overline{\mathbf Q}_\ell}
 =\{x\in\overline{\mathbf Q}_\ell:|x|_\ell<1\}$.

\begin{theorem}[K17-S17-JOINT]\label{thm:step17-joint}
Assume $\mathrm{GD}(\Omega,K,\ell,e,\psi)$. If the source Selmer
length $V_{\mathrm{len}}(\Omega,c,e,\psi)$ is nonzero, there is a
single representation $\rho\in\mathcal A$ such that
\begin{equation}\label{eq:joint-global-congruence}
 \operatorname{tr}\rho(s)-t(s)
       \in\mathfrak m_{\overline{\mathbf Q}_\ell}
       \qquad\text{for every }s\in G_K.
\end{equation}
Consequently, if every $\rho\in\mathcal A$ has at least one $s\in G_K$
for which this difference is a unit, then $V_{\mathrm{len}}=0$.
The detecting element may depend on $\rho$; no cardinality restriction
on $\mathcal A$ is required.
\end{theorem}

\begin{proof}
We first take the coefficient DVR $O$ to be unramified over
$\mathbf Z_\ell$. Use exactly the closed inputs of
\texttt{SmallClassNumbers/CorGKOf.lean}, proof of \texttt{corGK\_of},
and the construction in
\texttt{DenseThresholdCases/FOneCharU.lean}, proof of
\texttt{oneChar\_unramU}. They give the datum $D$, marked element
$g_1$, local object $U$, complete local deformation ring $S.R$,
universal representation $S.r$, and graded model $M$. In particular,
the source supplies the finiteness of the weight-zero algebra,
\texttt{PtHyp}, and
\begin{equation}\label{eq:joint-vlen-model}
 V_{\mathrm{len}}\le M.\mathrm{len}.
\end{equation}
These objects and hypotheses do not depend on the later test element
$s$. The same proof's assertion \texttt{hadm} states that every
\texttt{TwoPt} point $f:S.R\to\overline{\mathbf Q}_\ell$ gives
\[
 \rho_f=D.\mathrm{rhoK}(\operatorname{push}(f,S.r))\in\mathcal A.
\]
It uses the strengthened predicates displayed above, including the
source's \texttt{WitA}, \texttt{Unl} and Type-II hypotheses.

There must be a \texttt{TwoPt} point. Otherwise apply
\texttt{UnitaryThreeGroups/ModelUpper.lean}, theorem
\texttt{model\_upper}, at the identity element. Its trace nonvanishing
hypothesis is vacuous because there are no such points. It gives
$M.\mathrm{len}<\infty$ and a finite covering set $T$ consisting of
\texttt{TwoPt} points, with
\[
 \frac{M.\mathrm{len}}{e(O/\mathbf Z_\ell)}
       \le 2\sum_{f\in T}v_\ell(\operatorname{trDiff}(1,f)).
\]
The set $T$ is empty. Since $e(O/\mathbf Z_\ell)=1$, the finite
length is zero, and \eqref{eq:joint-vlen-model} contradicts
$V_{\mathrm{len}}\ne0$. This is the complete empty-point argument in
\texttt{SmallClassNumbers/OnePoint.lean}, proof of \texttt{onePoint};
here we use the stronger \texttt{hadm} assertion from
\texttt{FOneCharU.lean}.

Choose such an $f$. Its local-point condition and the universal residual
condition give the following entrywise congruence, for every
$g\in G_{\mathbf Q}$ and every $a,b\in\{0,1,2\}$:
\begin{equation}\label{eq:joint-entry-congruence}
 \left|\bigl(\operatorname{push}(f,S.r)(g)\bigr)_{ab}
            -\bigl(r_0(D.\eta,D.\Psi_Q)(g)\bigr)_{ab}\right|_\ell<1.
\end{equation}
For completeness, the proof of \texttt{isPt\_push} in
\texttt{UnitaryThreeGroups/IsPtOf.lean} establishes this directly:
the difference of the corresponding entries in $S.R$ has residue zero
by \texttt{S.r.red} and \texttt{diagonal\_rep\_map\_entries}, hence
belongs to its maximal ideal. The condition \texttt{IsLocPt f} sends
that ideal into elements of norm less than one. That proof also shows
that all entries of the pushed representation have norm at most one.
Thus \eqref{eq:joint-entry-congruence} is an actual residual congruence,
not merely a characteristic-zero trace identity.

Restrict to $G_K$. The identity $D.\eta|_{G_K}=1$ identifies the matrix
part with $\rho_f$. The exact trace identity \texttt{htr} in
\texttt{FOneCharU.lean}, using
\texttt{trace\_reducible\_three\_dim} from
\texttt{TheoremsMuLinked/GaloisRepresentationTypes/OneChar.lean},
identifies the trace of the base representation with $t(s)$.
The ultrametric inequality applied to the three diagonal differences
in \eqref{eq:joint-entry-congruence} now proves
\eqref{eq:joint-global-congruence} for every $s$ and the same $\rho_f$.

Finally remove the coefficient restriction exactly as in the final
proof of \texttt{FOneCharU.lean}. The prime-to-$\ell$ order of $\psi$,
part of $\mathrm{GD}$, gives an unramified coefficient ring $O'$ and
a character $\psi'$ with the same values in
$\overline{\mathbf Q}_\ell$. The lemma \texttt{gd\_chHom} preserves
$\mathrm{GD}$. The complete theorem \texttt{coefIndep} in
\texttt{HeightAndVanishing/CoefIndep.lean} gives equivalence of
infinite lengths and equality of normalized lengths. These two facts,
and the positive finite ramification indices, imply equivalence of
vanishing: if the $O'$-length were zero, the $O$-length would be finite
with normalized length zero and hence would be zero. Therefore the
$O'$-length remains nonzero. Apply the preceding construction over
$O'$. The characters $\Psi$ and $\Psi'$ have identical values, and the
predicates in $\mathcal A$ do not involve the coefficient ring. This
gives precisely \eqref{eq:joint-global-congruence} over the original
datum. Its stated contrapositive follows immediately.
\end{proof}

\begin{proposition}[K17-S17-GROUP-RING]\label{prop:step17-group-ring}
Under the same hypotheses, take $s_1,\ldots,s_r\in G_K$ and
$b_1,\ldots,b_r\in\mathbf Z_\ell$. Define
\[
 d_b(\rho)=\sum_{j=1}^r b_j
       \bigl(\operatorname{tr}\rho(s_j)-t(s_j)\bigr).
\]
If $d_b(\rho)\ne0$ for every $\rho\in\mathcal A$, then the source
length is finite and there is a finite set $X\subset\mathcal A$,
without isomorphic repetitions, such that
\begin{equation}\label{eq:group-ring-bound}
 \frac{V_{\mathrm{len}}}{e(O/\mathbf Z_\ell)}
       \le 2\sum_{\rho\in X}v_\ell(d_b(\rho)).
\end{equation}
In particular, units for all $d_b(\rho)$ force $V_{\mathrm{len}}=0$.
This is a finite-support group-ring statement; no infinite series is
being evaluated.
\end{proposition}

\begin{proof}
We give the complete modification to \texttt{ModelUpper.lean}; all
remaining source arguments are unchanged. Work first over an
unramified coefficient ring. For each $s_j$, use the homogeneous
entry lifts \texttt{M.ent} of its image in $G_{\mathbf Q}$ and let
$a_j\in M.A_0$ be the sum of the three diagonal lifts minus the
coefficient-ring trace of the base representation. The source's
identities \texttt{ha0}, \texttt{hqa} and \texttt{hxa} give
\[
 a_j\in M.A_0,\qquad M.x(a_j)=0,\qquad
 f(M.q(a_j))=\operatorname{trDiff}(s_j,f).
\]
The finite sum $a=\sum_j b_j a_j$ has the same degree and base value,
and evaluation is the corresponding linear combination of trace
differences. This uses only the $O$-module structure of the homogeneous
degree-zero part; $\mathbf Z_\ell$ embeds in every coefficient ring.

The hypothesis \texttt{M5} in that proof follows exactly as before:
the source map \texttt{M.pts} carries every two-signed model point,
up to the grading torus, to a local \texttt{TwoPt} point. Degree-zero
elements are fixed by that torus, and the evaluated $a$ is nonzero by
the present hypothesis and \texttt{hadm}. Apply \texttt{LemmaM} to
this $a$. The remainder of \texttt{model\_upper} constructs finite
torus-orbit representatives by the same argument, compares the same normalized
length, and evaluates the selected homogeneous element; it uses no
further property of a single group element. It therefore gives the
same inequality with $\operatorname{trDiff}$ replaced by the finite
linear combination.

The isomorphism-injectivity argument \texttt{hinj} of
\texttt{FOneCharU.lean} also does not depend on the test element.
Combining it with \texttt{vlen\_le} yields
\eqref{eq:group-ring-bound}. To descend coefficients, use
\texttt{exists\_unram}, \texttt{gd\_chHom} and \texttt{coefIndep}
as in the preceding proof. The $b_j$ have the same images in both
coefficient rings, and all $t(s_j)$ agree, so the finite linear
combination, finiteness and normalized bound descend. For units all
valuations vanish, giving the last assertion.
\end{proof}

\begin{corollary}[Joint operator test]\label{cor:joint-hecke-test}
Suppose a finite list of integral Hecke modules covers the source
admissible representations simultaneously at a finite collection of
Frobenius elements. On a given module write the scaled operators as
$A_1,\ldots,A_r$ and their integral Eisenstein targets as
$u_1,\ldots,u_r$; assume the scaling factors are units at $\ell$.
If, on every module, the reductions of
\[
 A_1-u_1I,\ldots,A_r-u_rI
\]
have no common nonzero kernel vector over
$\overline{\mathbf F}_\ell$, then $V_{\mathrm{len}}=0$.
A sufficient certificate on a module is
\[
 \det\left(\sum_{j=1}^r c_j(A_j-u_jI)\right)
       \in O^\times
\]
for some integral $c_j$; the coefficients may differ between modules.
More generally, a unit maximal minor of the vertically stacked
matrices is sufficient.
\end{corollary}

\begin{proof}
If the length were nonzero, Theorem~\ref{thm:step17-joint} would give
one admissible $\rho$ satisfying every residual trace congruence.
Simultaneous coverage places its characteristic-zero Hecke eigensystem
in one of the indicated modules. Pass to a finite coefficient extension
containing the selected eigenvalues and a simultaneous eigenvector,
and scale that vector to be integral and primitive. The eigenvalues
of all $A_j-u_jI$ belong to the maximal ideal because their scaling
factors are units. Reducing the primitive vector gives a nonzero common
kernel vector, a contradiction. An invertible linear combination, or
a full-rank stacked matrix, rules out such a vector. This argument
requires simultaneous coverage of one global representation; separate
characteristic polynomials with no identified common module do not
supply that hypothesis.
\end{proof}

\paragraph{Relation with the finite hyperplane argument.}
For a finite family of $n<\ell$ integral residual trace functions, each
distinct from the Eisenstein function, choose finitely many detecting
elements. For each representation the coefficient vectors whose linear
combination vanishes form a proper $\mathbf F_\ell$-linear subspace;
in $\mathbf F_\ell^r$ each has at most $\ell^{r-1}$ elements.
Their union therefore misses a coefficient vector. Integer lifts give
unit linear combinations for the whole family, so
Proposition~\ref{prop:step17-group-ring} applies. Theorem
\ref{thm:step17-joint} avoids this size restriction by retaining the
actual local point rather than first choosing a test function.

\paragraph{The remaining uniform arithmetic question.}
For $\psi=1$, $N=2$ and even $e$, the same Selmer group is tested by
each imaginary quadratic datum satisfying $\mathrm{GD}$. A nonzero
group forces an everywhere residual Eisenstein eigensystem in an
admissible branch for every such datum. Finding a datum with no such
eigensystem is therefore sufficient for vanishing. The converse is not
asserted: the presence of an automorphic congruence need not produce a
nonzero Selmer class by this argument. Neither an existential length
bound nor isolated constant-term calculations alone exclude the
remaining simultaneous congruences.
